\begin{abstract}
本文以“分辨率折叠所致的观测非单射”为出发点：在有限窗模型中，折叠映射 $\Fold_m:\Omega_m\twoheadrightarrow X_m$ 诱导原像纤维 $\mathcal F_m(x)=\Fold_m^{-1}(x)$，从而在均匀微测度下得到粗结果概率
$
p_m(x)=\abs{\mathcal F_m(x)}/\abs{\Omega_m}.
$
证明该计数概率可在有效 Hilbert 表述下写成 Born 权重，从而无需将 Born 规则作为新增公理。并在此基础上给出有限分辨率读出的最小量子接口封装：状态为正算子，读出为效应分解（POVM），去相干为相位记录丢失所诱导的商化/通道化。
\end{abstract}

\paragraph{贡献点（最小主张清单）}
\begin{itemize}
  \item 以纤维计数定义粗概率 $p_m$，并与折叠语言中的纤维谱/熵量对齐。
  \item 构造 $\mathcal H_m=\CC^{X_m}$ 上的向量 $|\psi_m\rangle$，使得投影读出满足 $\Pr(x)=|\langle x|\psi_m\rangle|^2=p_m(x)$。
  \item 给出最小接口封装：有限读出抽象为 POVM；去相干对应相位信息的丢失/商化。
\end{itemize}

\paragraph{关键词}
折叠算子；原像纤维；计数测度；Born 权重；POVM；去相干；残差充分统计量。

